\chapter{Estado del Arte}

\section{Procesamiento de lenguaje natural en el dominio financiero}
El vocabulario financiero presenta ambigüedades semánticas que invalidan los diccionarios y modelos de propósito general. Como demostraron Loughran y McDonald~\cite{loughran2011liability}, términos habituales en informes corporativos (como \textit{liability}, \textit{tax} o \textit{exposure}) se clasifican erróneamente como negativos en léxicos genéricos de sentimiento.

Para superar esta limitación, la literatura reciente ha convergido hacia modelos basados en arquitectura \textit{Transformer} preentrenados sobre corpus económicos y regulatorios, destacando \textbf{FinBERT}~\cite{yang2020finbert}. Estudios posteriores como el de Lopez-Lira y Tang~\cite{lopezlira2023chatgpt} han demostrado asimismo el potencial de los modelos de lenguaje de gran escala (LLM) para anticipar retornos diarios a partir de titulares financieros, planteando un equilibrio clave entre precisión semántica, latencia y coste computacional en producción.

\section{Fusión multimodal y rigor temporal en finanzas}
La combinación de series temporales de precios con flujos textuales dispersos fue formalizada en trabajos de referencia como StockNet~\cite{xu2018stock}. Sin embargo, el éxito práctico de estos enfoques depende del rigor en la validación temporal para evitar fugas de información (\textit{data leakage}) y correlaciones seriales espurias~\cite{lopezdeprado2018advances, jansen2020ml}.

\section{Regresión cuantílica y explicabilidad (XAI)}
\textit{[Completar: fundamentar la transición desde clasificación binaria / regresión a la media hacia regresión cuantílica con Pinball Loss para modelar volatilidad condicional e incertidumbre, así como el uso de TreeSHAP para auditar la atribución de variables textuales frente a indicadores técnicos.]}

\section{Brecha identificada y contribución del TFM}
\textit{[Completar: resumir en tabla comparativa qué aporta este TFM frente a enfoques habituales (Gatekeeper Few-Shot + SEC 8-K + Anexos EX-* + Point-in-Time NYSE + Regresión Cuantílica al 80\,\% + arquitectura Always Free).]}
